\documentclass[10pt]{article}
%\pdfoutput=1 
\usepackage{../tex/NotesTeX}
\usepackage{CTEX}

%\usepackage{showframe}

\title{\begin{center}{\Huge \textit{一份不知所云的笔记}}\\{{\itshape （我随便乱写的）}}\end{center}}
\author{Shunz Dai\footnote{\href{https://ShunzDai.github.io/}{\textit{My Personal Blog}}}}


\affiliation{
Physics and Space Science College,China West Normal University
}

\emailAdd{196883@outlook.com}

\begin{document}

	\maketitle
	\flushbottom
	\newpage
	\pagestyle{fancynotes}
	\part{引言}	

	我从2019年9月就开始写这份笔记了，但写到今天也只有寥寥数页。原本的笔记也有这样一个引言部分，主要讲述了我整个大学四年的艰难岁月，但现在看来，那些文字只不过是自己感动自己罢了\footnote{我仍然保留了那些文字，放在附录\ref{sec:Z}中。}。现在，我重写了这份引言，用尽量简短的语言来阐述这份笔记做了些什么工作。

	我首先给出我这篇笔记的定位，那就是，这篇笔记是写给我自己看的，是对我四年以来的工作的总结。

	这份笔记中，我们将通篇使用Einstein求和约定，并始终保持上下指标的差异，这种差异原则上与度量张量升降指标无关，是一种源自事物本身的“对偶”属性。至于“为什么物理量甲是上指标，物理量乙是下指标”这样的问题，我给出这样的解释：一是你首先要约定一个称为“上指标”的东西，然后才能找到与“上”对立的那个“下”；二是为了在引入度量张量后，各个物理量的指标位置关系仍然成立。这个解释很抽象，仿佛是“只可意会不可言传”。事实上也确实是这样，我们应当在学习中不断思考“对偶”的意义，多加计算方能修成正果。

	这份笔记受微分几何的影响颇深，而微分几何是我大学四年以来尤其喜爱的数学工具。相比于线性代数，微分几何为代数问题赋予了几何的看法，这使得我们能够从不同方向去看同一事物，像盲人摸象一般拼凑出研究对象的各种性质；相比于微积分（当然，我指的是高等数学中的微积分，而不是数学分析中的微积分），微分几何更注重“基本结构”，它没有繁多的运算技巧，不必将问题拘谨在冗长的计算中。
	
	我时常向朋友举例说明线性代数、微积分的研究对象与微分几何的研究对象的差异，这一差异正如这两个问题间的差异：第一个问题是“305497855加406804039等于多少”，第二个问题是“1加2为什么等于2加1”。对于第一个问题，我认为稍稍用点心便能计算出结果，是只有难度没有深度的问题；而对于第二个问题，我的朋友往往不能理解，他们通常会笑着问我“这是什么问题，那你说为什么1加2等于2加1”，这时候我就会回答“因为加法满足交换律”。我们在小学数学课上就已经学过“交换律”“结合律”“分配律”之类的东西，当我们在大学听到“1加2为什么等于2加1”这种问题时，按理说我们应当立即给出答案才对，但事实是我们往往答不出来。其原因就是我们从小到大受过的数学训练的目的就不是为了记住这些基本结构，而是为了将数学作为一种“工具”进行应用。所以，第二个问题是既有难度也有深度的问题。如果你现在产生了“1加2等于2加1用到了加法交换律，这有什么难度”这样的困惑，那是因为你没懂我想表达的意思，你应当回炉重造。
	
	在繁琐冗长的计算训练中，我们逐步忘却了那些基本结构的作用。当然，这些论断不意味着我们在学微分几何时不需要任何实际计算、应用，只是说，我们在学习微分几何时更应该将注意力放在基本结构上，而简单的计算能够加深理解，我们不必将精力放在如同上述第一个问题的复杂计算上。我们将从微分几何的基本结构出发，一步步搭建出整座大厦。

	\newpage

	\part{分析力学}\label{Part:Analytical}
\section{Lagrange力学}\label{sec:Lagrange}
	\begin{margintable}\vspace{1.4in}\footnotesize
		\begin{tabularx}{\marginparwidth}{|X}
		Section~\ref{sec:Lagrange}. Lagrange力学\\
		Section~\ref{sec:Hamilton}. Hamilton力学\\
		
		\end{tabularx}
	\end{margintable}

	\newpage
\section{Hamilton力学}\label{sec:Hamilton}
	Lagrange力学描述的是位形流形上的力学规律，是比Newton力学更深刻的分析工具。然而，Lagrange力学也有不足：Euler-Lagrange方程是一个二阶方程，它在大多数情况下不易求解。如果能想办法让描述力学系统的方程的阶数降下来，则在求解运动规律时就方便多了。Hamilton力学由此而生。Hamilton力学是辛流形上的力学规律，其参数为广义坐标$x^\mu$以及广义动量$p_\mu$。Lagrange力学中有Lagrange函数$L(x^\mu;{\dot{x}}^\mu)$，对应的，Hamilton力学中也有Hamilton函数，我们将其表示为$H(x^\mu;p_\mu)$。由Hamilton函数逐步搭建的力学体系能提供丰富的信息，事实上，由Heisenberg建立的矩阵力学就来源于经典Hamilton力学的量子化。
\subsection{Hamilton方程}
	对Lagrange函数$L(x^\mu;{\dot{x}}^\mu)$变分，得到
	\begin{equation}\label{eq:delta L}
	\delta L=\frac{\partial L}{\partial x^\mu}\delta x^\mu+\frac{\partial L}{\partial {\dot{x}}^\mu}\delta {\dot{x}}^\mu,
	\end{equation}
	Lagrange函数具有能量的量纲，于是Lagrange函数对广义速度的偏导数具有动量的量纲.不妨将这个具有动量量纲的量称为“广义动量”.
	\begin{definition}
	广义动量$p_\mu$定义为
	\begin{equation}
		p_\mu:=\frac{\partial L}{\partial {\dot{x}}^\mu}.
	\end{equation}
	\end{definition}
	于是(\ref{eq:delta L})可以改写成
	\begin{equation}\label{eq:delta L'}
	\delta L=\frac{\partial L}{\partial x^\mu}\delta x^\mu+p_\mu\delta {\dot{x}}^\mu;
	\end{equation}
	再根据Euler-Lagrange方程得到
	\begin{equation*}
	\frac{\partial L}{\partial x^\mu}=\frac{d}{dt}\frac{\partial L}{\partial {\dot{x}}^\mu}={\dot{p}}_\mu,
	\end{equation*}
	所以(\ref{eq:delta L'})变成
	\begin{equation}\label{eq:delta L''}
	\delta L={\dot{p}}_\mu\delta x^\mu+p_\mu\delta {\dot{x}}^\mu.
	\end{equation}
	为了得到某个以$p_\mu$为变量的变分，我们对(\ref{eq:delta L''})进行Legendre变换得到
	\begin{equation*}
	\delta L={\dot{p}}_\mu\delta x^\mu+\delta(p_\mu{\dot{x}}^\mu)-{\dot{x}}^\mu\delta p_\mu,
	\end{equation*}
	移项得到
	\begin{equation}\label{eq:delta H}
	\delta(p_\mu{\dot{x}}^\mu-L)=-{\dot{p}}_\mu\delta x^\mu+{\dot{x}}^\mu\delta p_\mu.
	\end{equation}
	所以,若
	\begin{definition}
	定义Hamilton函数为
	\begin{equation}
		H:=p_\mu{\dot{x}}^\mu-L,
	\end{equation}
	\end{definition}
	则$H$就是以$x^\mu$和$p_\mu$为变量的函数.于是我们可以建立如下原理
	\begin{theorem}
	Hamilton方程组是指如下两个方程
	\begin{equation}
		\begin{split}\label{eq:H}
			\frac{\partial H}{\partial x^\mu}&=-{\dot{p}}_\mu;\\
			\frac{\partial H}{\partial p_\mu}&={\dot{x}}^\mu.
		\end{split}
	\end{equation}
	\end{theorem}
	可见，Hamilton力学研究的是1阶方程，但方程的数量比Lagrange力学增加了一倍.

\subsection{Poisson括号}
	设任意力学量$A$是$x^\mu$与$p_\mu$的函数$A(x^\mu(t);p_\mu(t))$，我们来研究$A$随时间变化的规律.由于
	\begin{equation}\label{eq:Poisson}
	\frac{dA}{dt}=\frac{\partial A}{\partial x^\mu}\frac{d x^\mu}{dt}+\frac{\partial A}{\partial p_\mu}\frac{dp_\mu}{dt},
	\end{equation}
	考虑到Hamilton方程(\ref{eq:H})，带入(\ref{eq:Poisson})得到
	\begin{equation}\label{eq:Poisson'}
	\frac{dA}{dt}=\frac{\partial A}{\partial x^\mu}\frac{\partial H}{\partial p_\mu}-\frac{\partial H}{\partial x^\mu}\frac{\partial A}{\partial p_\mu}.
	\end{equation}
	(\ref{eq:Poisson'})存在着一种反对称性，我们可以用一个双线性映射表示之.
	\begin{definition}
	以$x^\mu$与$p_\mu$为变量的力学量$X$与$Y$的Poisson括号为
	\begin{equation}\label{eq:PB}
		\{X,Y\}:=\frac{\partial X}{\partial x^\mu}\frac{\partial Y}{\partial p_\mu}-\frac{\partial Y}{\partial x^\mu}\frac{\partial X}{\partial p_\mu}，
	\end{equation}
	\end{definition}
	\begin{lemma}
		Poisson括号满足
		\begin{enumerate}
			\item 反对称性：$\{X,Y\}=-\{Y,X\}$;
			\item Jacobi恒等式：$\{\{X,Y\},Z\}+\{\{Z,X\}Y\}+\{\{Y,Z\},X\}=0$.
		\end{enumerate}
	\end{lemma}
	于是(\ref{eq:Poisson'})可用Poisson括号表示为
	\begin{equation}
	\frac{dA}{dt}=\{A,H\}.
	\end{equation}
	事实上Poisson括号构成一个\textbf{Lie代数}，我们将在\ref{sec:LieAlgebra}中进一步论述Lie代数的性质.
	
	Hamilton方程也可以修改成Poisson括号的形式
	\begin{equation}
	\begin{split}
		\{p_\mu,H\}&={\dot{p}_\mu};\\
		\{x^\mu,H\}&={\dot{x}}^\mu,
	\end{split}
	\end{equation}
	而$x^\mu$和$p_\mu$的Poisson括号为
	\begin{equation}
	\{x^\mu,p_\nu\}=\frac{\partial x^\mu}{\partial x^\rho}\frac{\partial p_\nu}{\partial p_\rho}-\frac{\partial p_\nu}{\partial x^\rho}\frac{\partial x^\mu}{\partial p_\rho}=\delta^\mu_\rho\delta^\rho_\nu=\delta^\mu_\nu,
	\end{equation}
	上式利用了广义坐标与广义动量相互独立的属性.回到$A$随时间变化的规律问题上，我们得到如下引理.
	\begin{lemma}
	若力学量$A(x^\mu;p_\mu)$与Hamiltonian的Poisson括号为零，即
	\begin{equation}
		\{A,H\}=0,
	\end{equation}
	则物理量$A$是一个\textbf{守恒荷}.
	\end{lemma}
	一个适当的例子能帮助我们理解守恒荷的物理意义.
	\begin{example}
	设局部坐标下某自由粒子的Lagrangian为
	$$L=\frac{1}{2}g_{\mu\nu}\dot{x}^\mu\dot{x}^\nu,$$
	则广义动量
	\begin{eqnarray*}
		p_\rho&&=\frac{\partial L}{\partial {\dot{x}}^\rho}=\frac{1}{2}g_{\mu\nu}\left(\frac{\partial \dot{x}^\mu}{\partial {\dot{x}}^\rho}\dot{x}^\nu+\dot{x}^\mu\frac{\partial \dot{x}^\nu}{\partial {\dot{x}}^\rho}\right)\\
		&&=\frac{1}{2}g_{\mu\nu}\left(\delta^\mu_\rho\dot{x}^\nu+\dot{x}^\mu\delta^\nu_\rho\right)=g_{\rho\mu}\dot{x}^\mu.
	\end{eqnarray*}
	因此$\dot{x}^\mu\equiv g^{\mu\nu}p_\nu$，于是Hamiltonian为
	\begin{eqnarray*}
		H=p_\mu{\dot{x}}^\mu-L=g^{\mu\nu}p_\mu p_\nu-\frac{1}{2}g^{\mu\nu}p_\mu p_\nu=\frac{1}{2}g^{\mu\nu}p_\mu p_\nu	
	\end{eqnarray*}
	用Poisson括号可得
	\begin{eqnarray*}
		\{p_\mu,H\}=0;\\
		\{H,H\}=0.
	\end{eqnarray*}	
	所以此时广义动量和Hamiltonian守恒.
	\end{example}

	Poisson括号构成了一个函数空间上的Lie代数，事实上，Poisson括号被量子化后就得到了对易括号，对易括号是某个算子空间上的Lie代数.Poisson括号与对易括号之间有如下对应法则
	\begin{equation}
	\{A,B\}\rightarrow\frac{1}{i\hbar}[\hat{A},\hat{B}],
	\end{equation}
	其中$[\hat{A},\hat{B}]:=\hat{A}\hat{B}-\hat{B}\hat{A}$.

\subsection{Noether定理}
	\begin{theorem}
	设$\epsilon$是一个无穷小量.若无穷小坐标变换$x^\mu\rightarrow x'^\mu=x^\mu+\epsilon f^\mu(x)$保持Hamiltonian不变，则$Q\equiv p_\mu f^\mu(x)$是一个守恒荷.
	\end{theorem}
	\begin{proof}
	要证$Q$是守恒荷，即证$\dot{Q}$为零:
	\begin{equation*}
		\begin{split}
			\dot{Q}&\equiv\frac{dQ}{dt}=\{Q,H\}\\
			&=\frac{\partial Q}{\partial x^\rho}\frac{\partial H}{\partial p_\rho}-\frac{\partial H}{\partial x^\rho}\frac{\partial Q}{\partial p_\rho}\\
			&=p_\mu \frac{\partial f^\mu}{\partial x^\rho}\frac{\partial H}{\partial p_\rho}-f^\rho\frac{\partial H}{\partial x^\rho}.
		\end{split}
	\end{equation*}
	下面来计算Hamiltonian的变分.由于坐标变换保持Hamiltonian不变，则有$\delta H=0$，于是
	\begin{equation}\label{pr:delta H}
		0=\delta H=\frac{\partial H}{\partial x^\mu}\delta x^\mu+\frac{\partial H}{\partial p_\mu}\delta p_\mu,
	\end{equation}
	其中$\delta x^\mu$可直接由坐标变换得到：
	\begin{equation}\label{pr:delta x}
		\begin{split}
			x^\mu\rightarrow x'^\mu&=x^\mu+\epsilon f^\mu(x)\\
			\Rightarrow\delta x^\mu&=\epsilon f^\mu(x);
		\end{split}
	\end{equation}
	$\delta p_\mu$要从广义动量的定义出发得到：
	\begin{eqnarray*}
		&&p_\mu=\frac{\partial L}{\partial \dot{x}^\mu}\\
		\Rightarrow &&p'_\mu=\frac{\partial L}{\partial \dot{x}'^\mu}=\frac{\partial L}{\partial \dot{x}^\nu}\frac{\partial \dot{x}^\nu}{\partial \dot{x}'^\mu}=p_\nu\frac{\partial \dot{x}^\nu}{\partial \dot{x}'^\mu}，
	\end{eqnarray*}
	由坐标变换可推得
	\begin{eqnarray*}
		&&x'^\mu=x^\mu+\epsilon f^\mu(x)\\
		\Rightarrow&&\dot{x}'^\mu=\dot{x}^\mu+\epsilon \frac{\partial f^\mu}{\partial x^\rho}\dot{x}^\rho\\
		\Rightarrow&&\frac{\partial \dot{x}'^\mu}{\partial \dot{x}^\nu}=\delta^\mu_\nu+\epsilon \frac{\partial f^\mu}{\partial x^\rho}\delta^\rho_\nu=\delta^\mu_\nu+\epsilon \frac{\partial f^\mu}{\partial x^\nu}\\
		\Rightarrow&&\frac{\partial \dot{x}^\nu}{\partial \dot{x}'^\mu}=\delta^\nu_\mu-\epsilon \frac{\partial f^\nu}{\partial x^\mu},
	\end{eqnarray*}
	其中最后一步可用反证法证明
	\begin{eqnarray*}
		\frac{\partial \dot{x}'^\mu}{\partial \dot{x}^\nu}\frac{\partial \dot{x}^\nu}{\partial \dot{x}'^\rho}=\left(\delta^\mu_\nu+\epsilon \frac{\partial f^\mu}{\partial x^\nu}\right)\left(\delta^\nu_\mu-\epsilon \frac{\partial f^\nu}{\partial x^\mu}\right)=\delta^\mu_\rho+O(\epsilon^2),
	\end{eqnarray*}
	省略高阶无穷小只剩单位矩阵，说明左侧两个矩阵确实是互逆的.但我们尚未证明逆矩阵的唯一性，为此，我们假设
	\begin{eqnarray*}
		\frac{\partial \dot{x}^\nu}{\partial \dot{x}'^\mu}=\delta^\nu_\mu+\epsilon h^\nu_\mu,
	\end{eqnarray*}
	其中$h^\nu_\mu$是以坐标为变量的函数.于是
	\begin{eqnarray*}
		\delta^\mu_\rho&&=\frac{\partial \dot{x}'^\mu}{\partial \dot{x}^\nu}\frac{\partial \dot{x}^\nu}{\partial \dot{x}'^\rho}\\
		&&=\left(\delta^\mu_\nu+\epsilon \frac{\partial f^\mu}{\partial x^\nu}\right)\left(\delta^\nu_\rho+\epsilon h^\nu_\rho\right)\\
		&&=\delta^\mu_\rho+\epsilon\left(\frac{\partial f^\mu}{\partial x^\rho}+h^\mu_\rho\right)+O(\epsilon^2)\\
		&&\Rightarrow\epsilon\left(\frac{\partial f^\mu}{\partial x^\rho}+h^\mu_\rho\right)=0,
	\end{eqnarray*}
	要使上式对于任意无穷小量$\epsilon$都成立，唯有	
	\begin{eqnarray*}
		h^\mu_\rho=-\frac{\partial f^\mu}{\partial x^\rho},
	\end{eqnarray*}
	这样就证明了解的唯一性.所以
	\begin{equation}\label{pr:delta p}
	\begin{split}
		p'_\mu&=p_\nu\frac{\partial \dot{x}^\nu}{\partial \dot{x}'^\mu}=p_\nu\left(\delta^\nu_\mu-\epsilon \frac{\partial f^\nu}{\partial x^\mu}\right)\\
		&=p_\mu-\epsilon p_\nu\frac{\partial f^\nu}{\partial x^\mu}\\
		\Rightarrow\delta p_\mu&=-\epsilon p_\nu\frac{\partial f^\nu}{\partial x^\mu}
	\end{split}
	\end{equation}	
	将(\ref{pr:delta x})和(\ref{pr:delta p})带入(\ref{pr:delta H})得到
	\begin{equation}\label{pr:delta H'}
	\begin{split}
		0=\delta H&=-\epsilon\left[p_\nu\frac{\partial f^\nu}{\partial x^\mu}\frac{\partial H}{\partial p_\mu}-f^\mu\frac{\partial H}{\partial x^\mu}\right]=-\epsilon\dot{Q}.
	\end{split}
	\end{equation}
	要使(\ref{pr:delta H'})对于任意的无穷小量$\epsilon$都成立，唯有$\dot{Q}=0$，得证.\\
	\end{proof}

	\begin{theorem}
	守恒荷$Q$是无穷小坐标变换的生成元.
	\end{theorem}

	力学量$A(x^\mu;p_\mu)$的变分为
	\begin{eqnarray*}
	\delta A=\frac{\partial A}{\partial x^\mu}\delta x^\mu+\frac{\partial A}{\partial p_\mu}\delta p_\mu,
	\end{eqnarray*}
	从上面对守恒荷的证明可知，$\delta x^\mu$和$\delta p_\mu$能够借助守恒荷$Q$表述为
	\begin{exercise}
		\begin{equation}\label{eq:generator}
			\begin{split}
				\delta x^\mu&=\epsilon f^\mu=\epsilon\{x^\mu,Q\};\\
				\delta p_\mu&=-\epsilon p_\nu\frac{\partial f^\nu}{\partial x^\mu}=\epsilon\{p_\mu,Q\}.
			\end{split}
		\end{equation}
	\end{exercise}
	
	(\ref{eq:generator})请读者自行证明.因此$\delta A$可以进一步写成
	\begin{eqnarray*}
	\delta A=\epsilon\left(\frac{\partial A}{\partial x^\mu}\{x^\mu,Q\}+\frac{\partial A}{\partial p_\mu}\{p_\mu,Q\}\right).
	\end{eqnarray*}
	在经典情形下，生成元构造尚不明显.但在Poisson括号量子化后，可以非常明显的展示出生成元构造.例如
	\begin{eqnarray*}
	x^\mu\rightarrow x'^\mu&&=x^\mu+\delta x^\mu\\
	&&=x^\mu+\epsilon[x^\mu,Q]\\
	&&=x^\mu-\epsilon Qx^\mu+\epsilon x^\mu Q\\
	&&=\left(1-\epsilon Q\right)x^\mu\left(1+\epsilon Q\right)\\
	&&=e^{-\epsilon Q}x^\mu e^{\epsilon Q},
	\end{eqnarray*}
	其中第四个等号省略了$O(\epsilon^2)$.
\newpage
	\part{广义相对论}\label{Part:GR}
\section{微分几何}\label{sec:DG}
	\begin{margintable}\vspace{1.4in}\footnotesize
		\begin{tabularx}{\marginparwidth}{|X}
		Section~\ref{sec:DG}. 微分几何\\
		Section~\ref{sec:structure}. 结构数学\\
		Section~\ref{sec:LieAlgebra}. Lie代数\\
		\end{tabularx}
	\end{margintable}
\subsection{微分流形}


\subsection{张量场与微分形式场}


\subsection{联络与曲率}
	对于任意标量场$\phi\in T_{(0,0)}(M)$，其外微分是一个张量场
	\begin{equation}
		\begin{split}
	 		d\phi=\frac{\partial\phi}{\partial x^\mu}dx^\mu=\frac{\partial\phi}{\partial {x}^\mu}\frac{\partial{x}^\mu}{\partial {x'}^\nu}d{x'}^\nu=\frac{\partial{\phi'}}{\partial {x'}^\nu}d{x'}^\nu.
		\end{split}
	\end{equation}
	这是说，一个自然的问题是：对于切矢量场或余切矢量场，其外微分是否仍然是张量场？我们以余切矢量场为例，计算其外微分在坐标变换下的性质.
	
	对于一个余切矢量场$\omega\in T_{(0,1)}(M)$，借助局部坐标系可将其表述为
	\begin{equation}\label{eq:domega}
		\begin{split}
			\omega=\omega_\mu dx^\mu=\omega_\mu\frac{\partial{x}^\mu}{\partial {x'}^\nu}d{x'}^\nu=\omega'_\nu d{x'}^\nu,
		\end{split}
	\end{equation}
	\ref{eq:domega}的全微分在坐标变换下满足
	\begin{equation}
		\begin{split}	
		d\omega &=d\omega'_\nu{\wedge}d{x'}^\nu\\
		&= d\left(\omega_\mu\frac{\partial{x}^\mu}{{\partial}{x'}^\nu}\right)\wedge d{x'}^\nu\\
		&=\frac{\partial{x}^\mu}{{\partial}{x'}^\nu}d\omega_\mu{\wedge}d{x'}^\nu+\omega_{\mu}d\left(\frac{\partial{x}^\mu}{\partial {x'}^\nu}\right){\wedge}d{x'}^\nu\\
		&=\left(\frac{\partial{\omega}_\mu}{{\partial}{x}^\sigma}\frac{\partial{x}^\sigma}{\partial {x'}^\rho}\frac{\partial{x}^\mu}{{\partial}{x'}^\nu}+{\omega}_\mu\frac{\partial^2{x}^\mu}{\partial {x'}^\rho{\partial}{x'}^\nu}\right)d{x'}^\rho{\wedge}d{x'}^\nu\\
		&=\frac{\partial\omega'_\nu}{{\partial}{x'}^\rho}d{x'}^\rho{\wedge}d{x'}^\nu.
		\end{split}
	\end{equation}
	注意到过程中产生了一个二阶偏微分项，与张量的变换规律对比可知，这一项破坏了张量的变换规律.但外代数的全反对称特性可以消除掉这一项带来的影响，即
	\begin{equation}
		\begin{split}
			\frac{\partial\omega'_\nu}{{\partial}{x'}^\rho}d{x'}^\rho{\wedge}d{x'}^\nu&=2\partial'_{[\rho}\omega'_{\nu]}{dx'}^{\rho}\otimes{dx'}^\nu\\
			&=2\left(\partial_{[\sigma}\omega_{\mu]}\frac{\partial{x}^\sigma}{\partial {x'}^\rho}\frac{\partial{x}^\mu}{{\partial}{x'}^\nu}+{\omega}_\mu\frac{\partial^2{x}^\mu}{\partial {x'}^{[\rho}{\partial}{x'}^{\nu]}}\right){dx'}^{\rho}\otimes{dx'}^\nu\\
		&=2\partial_{[\sigma}\omega_{\mu]}\frac{\partial{x}^\sigma}{\partial {x'}^\rho}\frac{\partial{x}^\mu}{{\partial}{x'}^\nu}{dx'}^{\rho}\otimes{dx'}^\nu,
		\end{split}
	\end{equation}
	所以$d\omega$实际上仍是张量场。但微分形式场说到底只是一种特殊的张量场。对于一般的张量场，我们无法通过张量场自身的全反对称性消除上述二阶偏微分项。并且，我们也没有定义过切矢量场的外微分。这些困难意味着外微分运算不能直接推广到张量空间$T_{(p,q)}(M)$上去，于是，我们希望构造出一种定义在张量空间上的新的微分运算，使得任意张量场在这种微分运算下满足张量的变换规律。为此，我们引入\textbf{联络}。

	\subsubsection{仿射联络}
	现代微分几何中的联络指的是一般矢量丛上的联络.这一节我们先讨论切丛$T_{(1,0)}(M)$上的联络（称为\textbf{仿射联络}），然后逐步构造出张量丛$T_{(p,q)}(M)$上的联络.下面，我们直接给出仿射联络的定义.
	\begin{definition}
		仿射联络是一个映射
		\begin{equation*}
			D:T_{(1,0)}(M)\rightarrow T_{(1,1)}(M),
		\end{equation*}
		它满足下列条件：
		\begin{enumerate}
			\item 对任意的$A,B\in T_{(1,0)}(M)$有 D(A+B)=D(A)+D(B);						
			\item 对任意的$A\in T_{(1,0)}(M)，\alpha\in T_{(0,0)}(M)$有$D(\alpha A)=d\alpha\otimes A+\alpha D(A)$.
		\end{enumerate}
	\end{definition}
	局部上，联络由一组1微分形式给出.我们先看自然标架场$\{\partial_\mu\}$的联络.命
	\begin{equation}
	D(\partial_\mu)=\omega^\rho_\mu \otimes\partial_\rho={\Gamma^\rho}_{\mu\nu}dx^\nu\otimes \partial_\rho,
	\end{equation} 
	其中$\omega^\rho_\mu={\Gamma^\rho}_{\mu\nu}dx^\nu$，${\Gamma^\rho}_{\mu\nu}$称为\textbf{联络系数}，它是局部坐标系中的光滑函数.以$\omega^\rho_\mu$作为矩阵$\omega$第$\rho$行第$\mu$列的元素，这样构造的矩阵$\omega$称为\textbf{联络方阵}.可见，任意两个联络间的差异完全体现在联络方阵上.现在，我们来看联络的变换规律.在标架变换下，由联络的定义立即得到			 
	\begin{equation}
		\begin{split}
		D(\partial'_\nu) &= D\left(\frac{\partial{x^\mu}}{\partial{x'}^\nu}{\partial}_\mu\right) \\
		&= d\left(\frac{\partial{x^\mu}}{\partial{x'}^\nu}\right)\otimes{\partial}_\mu+\frac{\partial{x^\mu}}{\partial{x'}^\nu}D\left({\partial}_\mu\right)\\
		&= \left[d\left(\frac{\partial{x^\rho}}{\partial{x'}^\nu}\right)+\frac{\partial{x^\mu}}{\partial{x'}^\nu}\omega^\rho_\mu\right]\otimes{\partial}_\rho\\
		&= \left[d\left(\frac{\partial{x^\rho}}{\partial{x'}^\nu}\right)\frac{\partial{{x'}^\sigma}}{\partial{x^\rho}}+\frac{\partial{x^\mu}}{\partial{x'}^\nu}\omega^\rho_\mu\frac{\partial{{x'}^\sigma}}{\partial{x^\rho}}\right]\otimes\partial'_\sigma\\
		&={\omega'}^\sigma_\nu\otimes\partial'_\sigma，
		\end{split}
	\end{equation}
	其中
	\begin{equation}\label{eq:connection}
	{\omega'}^\sigma_\nu=d\left(\frac{\partial{x^\rho}}{\partial{x'}^\nu}\right)\frac{\partial{{x'}^\sigma}}{\partial{x^\rho}}+\frac{\partial{x^\mu}}{\partial{x'}^\nu}\omega^\rho_\mu\frac{\partial{{x'}^\sigma}}{\partial{x^\rho}},
	\end{equation}
	这就是联络方阵在局部标架场改变时的变换公式.我们还可以进一步展开\ref{eq:connection}，得到
	\begin{equation}
	{{\Gamma'}^\sigma}_{\nu\lambda}=\frac{\partial^2{x^\rho}}{\partial{x'}^\lambda\partial{x'}^\nu}\frac{\partial{{x'}^\sigma}}{\partial{x^\rho}}+{\Gamma^\rho}_{\mu\alpha}\frac{\partial{x^\mu}}{\partial{x'}^\nu}\frac{\partial{x}^\alpha}{\partial{x'}^\lambda}\frac{\partial{{x'}^\sigma}}{\partial{x^\rho}},
	\end{equation}
	这正是通常意义下的联络变换公式.可见联络系数${\Gamma^\rho}_{\mu\nu}$并不是一个张量，但引入它能使$D(\partial_\mu)$遵循张量的变换规律.但稍作思考可以发现，若使用全反对称化的技巧，我们也能利用联络系数构造出一个张量.我们将在下一节给出具体讨论.

	切丛的联络在余切丛上诱导出一个联络（仍记为$D$），易看出它是从$T_{(0,1)}(M)$到$T_{(0,2)}(M)$的映射.
	\begin{remark}
		相比于外微分算子，可见余切丛上的诱导联络算子将值域从外形式丛的截面（元素即全反对称的张量场）扩展到了一般的张量丛的截面（元素即一般的张量场）.所以我们说联络是外微分算子在一般张量丛上的推广.
	\end{remark}
	现在现在我们来考察$\{\partial_\mu\}$的对偶标架$\{dx^\mu\}$的联络.$D(dx^\mu)$由下式确定
	$$d\left(\partial_\mu,dx^\nu\right)=\left(D(\partial_\mu),dx^\nu\right)+\left(\partial_\mu,D(dx^\nu)\right),$$
	我们设
	$$D(dx^\nu)={\omega^*}^\nu_\rho\otimes{dx}^\rho,$$
	于是
	$$\left(\partial_\mu,D(dx^\nu)\right)=\left(\partial_\mu,{\omega^*}^\nu_\rho\otimes{dx}^\rho\right)={\omega^*}^\nu_\rho\delta^\rho_\mu={\omega^*}^\nu_\mu.$$
	考虑到对偶标架满足
	$$\left(\partial_\mu,dx^\nu\right)=\delta^\nu_\mu,$$ 
	于是得到								 
	\begin{equation}
		\begin{split}
			{\omega^*}^\nu_\mu&=\left(\partial_\mu,D(dx^\nu)\right)=-\left(D(\partial_\mu),dx^\nu\right)\\
			&=-\left(\omega^\rho_\mu\otimes\partial_\rho,dx^\nu\right)=-{\omega}^\rho_\mu\delta^\nu_\rho=-{\omega}^\nu_\mu,
		\end{split}
	\end{equation}								 
	即
	$D(dx^\nu)=-{\omega}^\nu_\rho\otimes{dx}^\rho$.
	至此，我们就能计算任意$A\in T_{(1,0)}(M)$与$B\in T_{(0,1)}(M)$的联络了.借助局部坐标系，我们有：								 
	\begin{equation}
		\begin{split}
			D(A)&=D(A^\mu\partial_\mu)=dA^\mu\otimes\partial_\mu+A^{\mu}D(\partial_\mu)\\
			&=(\partial_{\rho}A^\mu+A^{\nu}{\Gamma^\mu}_{\nu\rho})dx^\rho{\otimes}\partial_\mu=\nabla_{\rho}A^{\mu}dx^\rho{\otimes}\partial_\mu;\\
			D(B)&=D(B_{\mu}{dx}^\mu)=dB_\mu\otimes{dx}^\mu+B_{\mu}D({dx}^\mu)\\
			&=(\partial_{\rho}B_\mu-B_\nu{\Gamma^\nu}_{\mu\rho})dx^\rho{\otimes}dx^\mu=\nabla_{\rho}B_{\mu}dx^\rho{\otimes}dx^\mu,
		\end{split}	
	\end{equation}
	其中									 
	\begin{equation}
		\begin{split}
			\nabla_{\rho}A^\mu&=\partial_{\rho}A^\mu+A^{\nu}{\Gamma^\mu}_{\nu\rho};\\
			\nabla_{\rho}B_{\mu}&=\partial_{\rho}B_\mu-B_\nu{\Gamma^\nu}_{\mu\rho}.
		\end{split}		
	\end{equation}								
	这正是通常意义下的\textbf{协变导数}运算.由此可见，协变导数是一种依赖于局部坐标系的分量表述.用类似的方法，我们能得到$T_{(p,q)}(M)$上的诱导联络.
	\begin{definition}
		$T_{(p,q)}(M)$
		设任意$S\in T_{(p,q)}(M)$，借助局部坐标系$\{x^\mu\}$，可将$S$的联络表述为
		\begin{equation}
			\begin{split}
				D(S)&=D({S^{\mu_1\dots\mu_p}}_{\nu_1\dots\nu_q}\partial_{\mu_1}\otimes\dots\otimes\partial_{\mu_p}\otimes{dx}^{\nu_1}\otimes\dots\otimes{dx}^{\nu_q})\\
				&=d({S^{\mu_1\dots\mu_p}}_{\nu_1\dots\nu_q})\otimes\partial_{\mu_1}\otimes\dots\otimes\partial_{\mu_p}\otimes{dx}^{\nu_1}\otimes\dots\otimes{dx}^{\nu_q}\\
				&+{S^{\mu_1\dots\mu_p}}_{\nu_1\dots\nu_q}D(\partial_{\mu_1})\otimes\dots\otimes\partial_{\mu_p}\otimes{dx}^{\nu_1}\otimes\dots\otimes{dx}^{\nu_q}\\
				&+\dots+{S^{\mu_1\dots\mu_p}}_{\nu_1\dots\nu_q}\partial_{\mu_1}\otimes\dots\otimes D({\partial_{\mu_p}})\otimes{dx}^{\nu_1}\otimes\dots\otimes{dx}^{\nu_q}\\
				&+{S^{\mu_1\dots\mu_p}}_{\nu_1\dots\nu_q}\partial_{\mu_1}\otimes\dots\otimes{\partial_{\mu_p}}\otimes D({dx}^{\nu_1})\otimes\dots\otimes{dx}^{\nu_q}\\
				&+\dots+{S^{\mu_1\dots\mu_p}}_{\nu_1\dots\nu_q}\partial_{\mu_1}\otimes\dots\otimes{\partial_{\mu_p}}\otimes{dx}^{\nu_1}\otimes\dots\otimes D({dx}^{\nu_q})\\
				&=(\partial_\rho{S^{\mu_1\dots\mu_p}}_{\nu_1\dots\nu_q}+{S^{\sigma\dots\mu_p}}_{\nu_1\dots\nu_q}{\Gamma^{\mu_1}}_{\sigma\rho}\\
				&+\dots+{S^{\mu_1\dots\sigma}}_{\nu_1\dots\nu_q}{\Gamma^{\mu_p}}_{\sigma\rho}-{S^{\mu_1\dots\mu_p}}_{\sigma\dots\nu_q}{\Gamma^{\sigma}}_{\nu_1\rho}\\
				&-\dots-{S^{\mu_1\dots\mu_p}}_{\nu_1\dots\sigma}{\Gamma^{\sigma}}_{\nu_q\rho})dx^\rho\otimes\partial_{\mu_1}\otimes\dots\otimes\partial_{\mu_p}\otimes{dx}^{\nu_1}\otimes\dots\otimes{dx}^{\nu_q}\\
				&=\nabla_{\rho}{S^{\mu_1\dots\mu_p}}_{\nu_1\dots\nu_q}dx^\rho\otimes\partial_{\mu_1}\otimes\dots\otimes\partial_{\mu_p}\otimes{dx}^{\nu_1}\otimes\dots\otimes{dx}^{\nu_q},
			\end{split}		
		\end{equation}
		其中
		\begin{equation}
			\begin{split}
				\nabla_{\rho}{S^{\mu_1\dots\mu_p}}_{\nu_1\dots\nu_q}&=\partial_\rho{S^{\mu_1\dots\mu_p}}_{\nu_1\dots\nu_q}+{S^{\sigma\dots\mu_p}}_{\nu_1\dots\nu_q}{\Gamma^{\mu_1}}_{\sigma\rho}\\
				&+\dots+{S^{\mu_1\dots\sigma}}_{\nu_1\dots\nu_q}{\Gamma^{\mu_p}}_{\sigma\rho}-{S^{\mu_1\dots\mu_p}}_{\sigma\dots\nu_q}{\Gamma^{\sigma}}_{\nu_1\rho}\\
				&-\dots-{S^{\mu_1\dots\mu_p}}_{\nu_1\dots\sigma}{\Gamma^{\sigma}}_{\nu_q\rho}.
			\end{split}		
		\end{equation}
	\end{definition}

\newpage

\part{量子场论}\label{Part:QFT}
\section{高等量子力学}\label{sec:QM}
	\begin{margintable}\vspace{1.4in}\footnotesize
		\begin{tabularx}{\marginparwidth}{|X}
		Section~\ref{sec:DG}. 量子化\\
		Section~\ref{sec:structure}. 路径积分\\
		Section~\ref{sec:LieAlgebra}. Lie代数\\
		\end{tabularx}
	\end{margintable}

	\section{规范场论}\label{sec:GFT}
	

\end{document}